\documentclass[a4paper,11pt]{article}
\usepackage{amsmath,amsthm,amsfonts,amssymb,amscd,amstext,vmargin,graphics,graphicx,tabularx,multicol} 
\usepackage[francais]{babel}
\usepackage[utf8]{inputenc}  
\usepackage[T1]{fontenc} 
\usepackage{pstricks-add,tikz,tkz-tab,variations}
\usepackage[autolanguage,np]{numprint} 
\usepackage{xlop}


\setmarginsrb{1.5cm}{0.5cm}{1cm}{0.5cm}{0cm}{0cm}{0cm}{0cm} %Gauche, haut, droite, haut
\newcounter{numexo}
\newcommand{\exo}[1]{\stepcounter{numexo}\noindent{\bf Exercice~\thenumexo} : \marginpar{\hfill /#1}}
\reversemarginpar


\newcounter{enumtabi}
\newcounter{enumtaba}
\newcommand{\q}{\stepcounter{enumtabi} \theenumtabi.  }
\newcommand{\qa}{\stepcounter{enumtaba} (\alph{enumtaba}) }
\newcommand{\initq}{\setcounter{enumtabi}{0}}
\newcommand{\initqa}{\setcounter{enumtaba}{0}}

\newcommand{\be}{\begin{enumerate}}
\newcommand{\ee}{\end{enumerate}}
\newcommand{\bi}{\begin{itemize}}
\newcommand{\ei}{\end{itemize}}
\newcommand{\bp}{\begin{pspicture*}}
\newcommand{\ep}{\end{pspicture*}}
\newcommand{\bt}{\begin{tabular}}
\newcommand{\et}{\end{tabular}}
\renewcommand{\tabularxcolumn}[1]{>{\centering}m{#1}} %(colonne m{} centrée, au lieu de p par défault) 
\newcommand{\tnl}{\tabularnewline}

\newcommand{\bmul}[1]{\begin{multicols}{#1}}
\newcommand{\emul}{\end{multicols}}

\newcommand{\trait}{\noindent \rule{\linewidth}{0.2mm}}
\newcommand{\hs}[1]{\hspace{#1}}
\newcommand{\vs}[1]{\vspace{#1}}

\newcommand{\N}{\mathbb{N}}
\newcommand{\Z}{\mathbb{Z}}
\newcommand{\R}{\mathbb{R}}
\newcommand{\C}{\mathbb{C}}
\newcommand{\Dcal}{\mathcal{D}}
\newcommand{\Ccal}{\mathcal{C}}
\newcommand{\mc}{\mathcal}

\newcommand{\vect}[1]{\overrightarrow{#1}}
\newcommand{\ds}{\displaystyle}
\newcommand{\eq}{\quad \Leftrightarrow \quad}
\newcommand{\vecti}{\vec{\imath}}
\newcommand{\vectj}{\vec{\jmath}}
\newcommand{\Oij}{(O;\vec{\imath}, \vec{\jmath})}
\newcommand{\OIJ}{(O;I,J)}


\newcommand{\reponse}[1][1]{%
\multido{}{#1}{\makebox[\linewidth]{\rule[0pt]{0pt}{20pt}\dotfill}
}}

\newcommand{\titre}[5] 
% #1: titre #2: haut gauche #3: bas gauche #4: haut droite #5: bas droite
{
\noindent #2 \hfill #4 \\
#3 \hfill #5

\vspace{-1.6cm}

\begin{center}\rule{6cm}{0.5mm}\end{center}
\vspace{0.2cm}
\begin{center}{\large{\textbf{#1}}}\end{center}
\begin{center}\rule{6cm}{0.5mm}\end{center}
}



\begin{document}
\pagestyle{empty}
\titre{Correction du contrôle sur les divisions euclidiennes }{Nom :}{Prénom :}{Classe}{Date}




\vspace*{0.35cm}


\exo{3} \includegraphics[scale=0.4]{trefle.eps}COURS

\q  \textbf{Voir la leçon !}\\



\q Compléter les phrases avec le vocabulaire suivant : un multiple,  un diviseur ou divisible.      
 
 \bmul{2}
  35 est \textbf{un multiple}  7. \\
  
  1448 est \textbf{ divisible par}  2 et 4.\\ 
   
  \columnbreak
  
  1 450 est \textbf{un multiple de}  10.\\
  
  8 est \textbf{un diviseur de}   56 .\\
  
                                         
\emul
  
 

              \vspace*{0.25cm}


\exo{3} Poser et effectuer les divisions euclidiennes sur votre copie. \\
Pour chaque division ci-dessous, \textbf{écrire l'égalité} vue dans le cours.\\

\initqa
\hspace*{1.5cm} \qa 2716 $\div$ 3            \hspace*{6cm}                    \qa 7 869 $\div$ 25       \\

\hspace*{1.5cm}  $2 716 = 3 \times 905 + 1$          \hspace*{6cm}                     $7 869 = 25 \times 314 + 19$       \\


              \vspace*{0.5cm}
              
 \exo{1.5}
  Lorsque l'on effectue la division euclidienne  d'un nombre par 13, le quotient est 9 et le reste est 7.\\
  Quel est ce nombre ? \textit{(Justifier votre réponse.)}\\
  
  \textbf{$13 \times 9 +7=117+7=124$ Le nombre cherché est 124.}
              \vspace*{0.5cm}

\exo{3.5}
\initq

\q  \textbf{Voir la leçon !}\\


\q Parmi les nombres suivants : 
 \hspace*{0.25cm}; \hspace*{0.25cm} 233 \hspace*{0.25cm} ; \hspace*{0.25cm} 2115 \hspace*{0.25cm} ;\hspace*{0.25cm} 2523 \hspace*{0.25cm}; \hspace*{0.25cm} 210  \hspace*{0.25cm} ; \hspace*{0.25cm} 468 \hspace*{0.25cm} ;\hspace*{0.25cm} 57\\


\noindent \initqa \qa	Quel(s) sont ceux qui sont divisibles par 2 ?\\
\textbf{Ceux qui sont divisibles par 2 sont 210 et 468.}\\

\qa	Quel(s) sont ceux qui sont divisibles par 3 ?\\
\textbf{Ceux qui sont divisibles par 3 sont: 2 115 ; 2 523 ; 210 ; 468 ; 57  }\\

\qa	Quel(s) sont ceux qui sont divisibles par 4 ?\\
\textbf{Celui qui est divisible par 4 est 468.}\\

\qa	Quel(s) sont ceux qui sont divisibles par 9 ?\\
\textbf{Ceux qui sont divisibles par 9 sont : 2 115 et 468.}

              \vspace*{0.75cm}
              
  \exo{2} "Le ticket gagnant est un nombre qui est à la fois multiple de 4 et multiple de 9."
  
  \includegraphics[scale=1.3]{ticketgagnant.eps} 
  
  $\rightarrow$ \textbf{Quel est le ticket gagnant ? Justifier votre réponse.           }\\
  
  \textbf{C'est le ticket 9 756. 9 + 7 + 5 + 6 = 27 et 27 est divisible par 9 donc 9 756 est divisible par 9.}\\
  \textbf{De plus $56 = 4 \times 14$ donc 9 756 est divisible par 4 aussi.}\\
              \newpage

\exo{2} Pour le C.D.I. du collège, la documentaliste reçoit 370 livres qu'elle doit ranger sur des étagères. Elle ne peut transporter que 13 livres à la fois. \\

\initq \q Combien de voyages devra-t-elle faire au minimum ? \textit{(Justifier votre réponse.)}\\

\textbf{$370= 13 \times 28 +6$ donc elle devra faire 29 voyages minimum pour ranger tous les livres.}\\

\q Combien de livres transportera-t-elle au dernier voyage ?\\

\textbf{Il restera 6 livres au dernier voyage.}

\vspace*{0.75cm}

\exo{3} Dans un collège, 163 élèves sont inscrits à
l'UNSS. Le responsable veut acheter un maillot
pour chacun des inscrits. Les maillots sont vendus
par lot de 14.\\

\initq \q Combien de lots doit-il acheter ?\\

\textbf{$163 = 11 \times 14 + 9$ Il doit donc acheter 12 lots de maillots.}\\

\q Combien de maillots ne seront pas distribués ?\\

\textbf{$12 \times 14 =168$ et $168 - 163 = 5$ Il restera donc 5 maillots non utilisés.}\\

              \vspace*{0.75cm}



\end{document}
